% ============================================================================
%  SECTION 10 — Open-source implementations
% ============================================================================
\begin{frame}[t]{}
  \vfill\begin{center}{\Huge\color{YaleBlue} Open-Source Implementations}
  \end{center}\vfill
\end{frame}

% ---- how the seminar code works: file-by-file rundown -----------------------
\begin{frame}[t]{Under the Hood, What Each File Computes (1/2: the core)}
  \footnotesize
  The \texttt{gr} package is layered: a shared \textbf{continuum-mechanics core}
  (read, don't edit), then one file per G\&R theory.
  \vspace{0.2em}
  \renewcommand{\arraystretch}{1.4}
  \begin{tabular}{@{}p{0.21\textwidth}p{0.42\textwidth}p{0.30\textwidth}@{}}
    \toprule
    \textbf{File} & \textbf{Computes} & \textbf{Continuum mechanics / G\&R} \\
    \midrule
    \texttt{mechanics.py} & strain energy $W(\lambda_e)$, 2nd PK stress
      $S=\dd W/\dd E_e$, Cauchy $\sigma=\lambda_e^2 S$ (neo-Hookean, Fung) &
      finite-strain \textbf{hyperelasticity} of each constituent, in reference
      measures \\
    \texttt{parameters.py} & \texttt{Model}/\texttt{Constituent}
      ($\phi_0,G^{\alpha},k_d,K_\sigma$); set-point
      $\sigma_h^{\alpha}=\sigma^{\alpha}(G^{\alpha})$; \texttt{Insult} &
      \textbf{deposition prestretch} \& homeostatic \textbf{set-points}; the two
      insults \\
    \texttt{geometry.py} & required Laplace stress
      $\bar\sigma_h\,\gamma\,\lambda^2/m$; solves
      $\bar\sigma(\lambda)=\sigma_{\rm req}$ for $\lambda$ &
      \textbf{mechanical equilibrium}; no root $\Rightarrow$ runaway \\
    \texttt{history.py} & \texttt{Result}: $t,\lambda,M/M_0,\bar\sigma,\dots$ &
      shared time-history output \\
    \bottomrule
  \end{tabular}
\end{frame}

\begin{frame}[t]{Under the Hood, What Each File Computes (2/2: the theories)}
  \footnotesize
  \renewcommand{\arraystretch}{1.3}
  \begin{tabular}{@{}p{0.25\textwidth}p{0.41\textwidth}p{0.27\textwidth}@{}}
    \toprule
    \textbf{File} & \textbf{Computes} & \textbf{G\&R theory / role} \\
    \midrule
    \texttt{kinematic\_growth.py} & split $\bm{F}=\bm{F}_e\bm{F}_g$, scalar
      $\theta$; $\dot\theta=k_g\theta(\bar\sigma/\bar\sigma_h-1)$ &
      Theory 1: kinematic growth \\
    \texttt{constrained\_mixture.py} & heredity integrals over cohorts; mass
      balance + mixture stress ($O(N^2)$) & Theory 2: full CMM \\
    \texttt{homogenized\_cmm.py} & two ODEs per constituent: $\dot M^{\alpha}$,
      $\dot\lambda_n^{\alpha}$ ($O(N)$) & Theory 3: homogenized CMM \\
    \texttt{equilibrated\_cmm.py} & rates $\to 0$: one algebraic solve for
      $\lambda^*$ & Theory 4: equilibrated CMM \\
    \texttt{stability.py} & \texttt{adapts} (does a root exist?),
      \texttt{critical\_elastin\_loss} & stability $=$ \textbf{existence} of an
      equilibrium \\
    \texttt{scenario.py}, \texttt{run.py}, \texttt{configs/} & build
      model/geometry/insult from YAML; run a study; sweep & the exercise driver
      (input files) \\
    \bottomrule
  \end{tabular}
  \vspace{0.3em}
  \begin{center}\footnotesize
  Common thread: each theory changes only \emph{how the supplied stress
  $\bar\sigma(\lambda)$ is built}, the same constituent laws and equilibrium
  solve are reused, so every difference you saw is the G\&R law alone.
  \end{center}
\end{frame}

\begin{frame}[t]{Research-Grade Open-Source Solvers}
  \small
  \begin{center}
  \renewcommand{\arraystretch}{1.4}
  \begin{tabular}{@{}p{0.16\textwidth}p{0.44\textwidth}p{0.30\textwidth}@{}}
    \toprule
    \textbf{Project} & \textbf{What it does} & \textbf{Stack / group} \\
    \midrule
    \textbf{\color{YaleBlue}svGrowth} & constrained-mixture G\&R of vessels
      (biaxial thin-wall); Stanford CBCL \\
    \textbf{\color{YaleBlue}svFSGe} & fast, strongly-coupled 3D
      fluid--solid--growth driver (the FSGe method) &
      Python + svFSI/svMultiPhysics; Yale (Pfaller lab) \\
    \textbf{\color{YaleBlue}FEMBE\_Plugin} & mechanobiologically equilibrated
      constrained-mixture material as a FEBio plugin &
      C++ / FEBio; Yale (Humphrey lab) \\
    \textbf{\color{YaleBlue}FEFSG\_Plugin} & fluid--solid--growth coupling as a
      FEBio plugin (thoracic-aorta hypertension \& aneurysm cases) &
      C++ / FEBio; Yale (Humphrey lab) \\
    \bottomrule
  \end{tabular}
  \end{center}
  \vspace{0.5em}
  {\footnotesize
  \texttt{github.com/StanfordCBCL/svGrowth} \quad
  \texttt{github.com/StanfordCBCL/svFSGe} \\
  \texttt{github.com/yale-humphrey-lab/FEMBE\_Plugin} \quad
  \texttt{github.com/yale-humphrey-lab/FEFSG\_Plugin}}
  \srccite{svGrowth builds on Humphrey, Rajagopal, Math Models Methods Appl Sci (2002); FEMBE on Latorre, Humphrey, Comput Methods Appl Mech Eng (2020)}
\end{frame}

